\chapter{Research and Development Activities}
\label{ch:activities}

This chapter reconstructs the project as a chronological scientific and
engineering investigation. Each section addresses one implementation phase in
turn, but the organising logic is not "what was built" so much as: what
question the phase was trying to answer, what the resulting evidence showed,
what limitation that evidence exposed, and what decision followed from it.
Provenance labels for each phase are given only in footnotes and figure/table
captions, where they matter for traceability; the prose refers to phases by
their scientific role. The shared methodological vocabulary — specimen-level
grouping, the three evaluation regimes, the representation-strategy taxonomy,
and the distinction between observed, model-derived, and derived-screening
quantities — was established in Chapter~\ref{ch:framework} and is used here
without re-derivation.

\section{Exploratory Prototype}
\label{sec:phase0}

The project's first implementation was an informal exploratory prototype,
built to establish whether the available photographs contained any usable
signal at all rather than to produce a validated model.\footnote{Corresponds
to the earliest exploratory implementation referenced in Chapter~\ref{ch:dataset}'s
corrupted-image-handling table.} It computed a single fixed-threshold
colour mask per image, derived a handful of additional descriptors by parsing
the specimen identifier as a string (rather than reading the structured
metadata table), and trained a RandomForest to predict next-week visible rust
percentage. A separate script then applied this one-step regressor
autoregressively — repeatedly feeding each week's prediction back in as the
next week's input — until a fixed 10\% rust threshold was crossed, and
reported the number of simulated steps as a "remaining useful life."

This should not be read as a trained or validated RUL model: it is a
heuristic simulation built on a single-step regressor, compounding that
regressor's error at every simulated step, with no mechanism to check the
simulated trajectory against an observed one. More consequentially, the
script that produces the prototype's "production" model fits on the entire
available dataset with no held-out partition at all, and no metrics from
this phase were saved anywhere in the repository — every quantitative claim
about this phase's performance would have to come from console output that
no longer exists. The phase is therefore useful only as a starting point: it
suggested that a colour-based signal was worth pursuing, while simultaneously
demonstrating that a benchmarked, leakage-aware, and reproducible pipeline
did not yet exist. That gap is what the next phase was built to close.

\section{Classical Machine-Learning Baseline}
\label{sec:phase1}

The second phase replaced the prototype with a benchmarked pipeline: a
consistent set of candidate regressors (ElasticNet, RandomForest, Gradient
Boosting, and XGBoost) evaluated under a specimen-grouped holdout followed by
grouped cross-validation on the training partition, with metrics and feature
importances persisted for every candidate. Rather than posing one prediction
problem, the phase implemented three distinct tasks over the same
underlying data, distinguished here because they have different scientific
meanings and must not be merged.

\textbf{Current visible-corrosion estimation} predicted the current-week peak
rust percentage from a combination of handcrafted image descriptors and
specimen/exposure metadata. The selected RandomForest achieved a grouped
holdout MAE of 2.38 and $R^2 = 0.900$ — the strongest headline number produced
by this phase. Inspecting the model's own feature importances, however,
shows that a single handcrafted feature — the percentage of pixels falling
inside a fixed rust-colour threshold — accounts for over 90\% of the total
importance. Because this same colour threshold is also how the target label
itself is effectively derived, the result is close to a restatement of the
labelling rule rather than independent evidence that the model has learned a
generalisable visual concept of corrosion. The headline accuracy is real, but
it does not by itself establish that the model has captured visual
information beyond the threshold used to construct its own label.

\textbf{Progression modelling} predicted the same current-week rust target
from specimen/exposure metadata and observation week alone, with no image
input at all. This is a metadata- and week-conditioned regression, not a
sequence model: it does not consume a specimen's own history of prior
observations and does not forecast forward from one time step to the next.
The selected Gradient Boosting model reached a grouped holdout $R^2$ of only
0.517, substantially weaker than the image-assisted task above, and the
result should be read accordingly — as an estimate of how much predictive
information design and exposure variables alone carry for this target,
not as a temporal forecasting result.

\textbf{Threshold-timing estimation} predicted the number of weeks remaining
until peak rust crosses a fixed 20\% threshold, using the specimen's true
current corrosion level (not a model's own prediction of it) together with
metadata as inputs. The selected RandomForest reached a grouped holdout MAE
of 2.21 weeks. This is a derived timing quantity conditioned on an observed
covariate, not an instance of remaining-useful-life prediction in the sense
defined in Section~\ref{sec:proxy-rul}: it times a corrosion threshold, not a
structural failure or service-life endpoint, and it is trained on already
known rather than forecast inputs.

\begin{table}[htbp]
\centering
\caption{Representative grouped-holdout results, classical baseline phase.
Selected model and holdout metrics as saved by the phase's own evaluation
script.}
\label{tab:phase1-results}
\begin{tabular}{p{4.3cm}p{2.6cm}p{1.5cm}p{1.6cm}p{1.5cm}}
\toprule
Task & Selected model & MAE & RMSE & $R^2$ \\
\midrule
Current corrosion (image + metadata) & RandomForest & 2.38 & 4.25 & 0.900 \\
Progression (metadata + week only) & Gradient Boosting & 6.71 & 9.33 & 0.517 \\
Threshold-timing (weeks to 20\%) & RandomForest & 2.21 wk & 3.09 wk & 0.782 \\
\bottomrule
\end{tabular}
\end{table}

Two further findings from this phase materially affect how its results
should be read. First, the specimen-grouped holdout used throughout this
phase was a genuine improvement over the exploratory prototype's absence of
any held-out evaluation, and is the same grouping principle formalised in
Section~\ref{sec:leakage}. Second, model selection within each task was
based on the holdout metric itself rather than on a separate validation
partition: for the threshold-timing task, an alternative candidate achieves
a better holdout RMSE (2.82 vs.\ 3.09) and $R^2$ (0.818 vs.\ 0.782) than the
selected RandomForest, which is chosen because its holdout MAE is
marginally lower (2.21 vs.\ 2.24). This is a historical model-selection
practice specific to this phase, not the methodological standard adopted
later in the project (Section~\ref{sec:principles}), and it means the
reported holdout numbers should be read as an optimistic upper bound on
this phase's true generalisation performance rather than as an untouched
final estimate. The pipeline's reproducibility is also limited in the
current repository state: re-running it against the presently stored data
workbook raises a schema mismatch, so the numbers above are the saved
record of a historical run rather than a result that can be freshly
regenerated without repair.

Taken together, the phase established that a benchmarked, specimen-grouped
pipeline could produce strong-looking results for current-corrosion
estimation, but left open whether that strength reflected a genuinely useful
image representation or simply a restatement of the handcrafted threshold
used to build the label. Testing whether a representation learned directly
from image content, rather than a fixed hand-designed threshold, could
resolve this question motivated the next phase.

\section{Deep Image-Embedding Experiment}
\label{sec:phase2}

The third phase replaced the handcrafted colour threshold with a
representation learned by a convolutional network: each image was passed
through a pretrained ImageNet ResNet-18 with its final classification layer
removed, used purely as a frozen feature extractor (no weights in the
backbone were updated), producing a 512-dimensional embedding. The
experiment targeted the same family of quantity as the previous phase's
strongest result — current-week peak rust percentage — so that the value of
a learned representation could be assessed against a comparable target,
computed here on the phase's own 792-row dataset basis, with the one
corrupted image's embedding replaced by the column-wise median.

Two configurations were compared: an image-only regressor consuming only the
512-dimensional embedding, and a multimodal regressor consuming the same
embedding concatenated with seven metadata fields. No metadata-only baseline
was trained in this phase, so the comparison that later became central to
the project's methodology (Section~\ref{sec:multimodal}) is only partially
available here.

The results contain a genuine reversal worth stating clearly.
Across training, the multimodal model reached a lower best validation MAE
than the image-only model (4.02 vs.\ 4.96). On the held-out test partition,
however, the ranking inverted: the image-only model achieved a lower test
MAE than the multimodal model (5.63 vs.\ 5.75). Figure~\ref{fig:phase2-fusion}
shows both comparisons side by side. The phase's own pipeline selected the
final deployed model by test MAE — the same partition used to report the
final metric — so the image-only model was deployed on the strength of a
test-set comparison that also served as the reported result, rather than
being confirmed by an independent partition after selection. This mirrors
the model-selection pattern already noted in Section~\ref{sec:phase1} and
should likewise be read as a historical practice specific to this phase,
not a result that establishes multimodal fusion as inferior in general.

\begin{figure}[htbp]
\centering
\includegraphics[width=0.85\textwidth]{fig_4_1_main2_image_only_vs_multimodal.png}
\caption{Deep-embedding phase: best validation MAE versus test MAE for the
image-only and multimodal (image + metadata) regressors, current-week peak
rust percentage. The multimodal model is ahead on validation but behind on
test; the deployed model was chosen by the test-partition comparison shown
on the right. Generated from the phase's own saved training history and
model-comparison tables.}
\label{fig:phase2-fusion}
\end{figure}

A second finding concerns where the deployed model's errors concentrate.
Grouping the saved test-set residuals by observation week shows a
systematic pattern: mean absolute error rises from 2.22 at week 10 or
earlier, to 6.62 across weeks 12–20, to 12.08 from week 23 onward, and the
average bias moves from a small positive value ($+0.85$) to a substantial
negative one ($-6.91$) over the same progression. Figure~\ref{fig:phase2-residuals}
shows this directly: the model increasingly underestimates severity as
corrosion advances, rather than erring in either direction at random.

\begin{figure}[htbp]
\centering
\includegraphics[width=0.75\textwidth]{fig_4_2_main2_residual_vs_week.png}
\caption{Deep-embedding phase, image-only model: residuals (predicted minus
observed peak rust percentage) on the held-out test partition, by
observation week and specimen series. Residuals drift increasingly negative
at later weeks, indicating systematic underestimation of severe, late-stage
corrosion. Reproduced from the phase's own saved diagnostic figure.}
\label{fig:phase2-residuals}
\end{figure}

These findings support a narrower conclusion than either "deep embeddings
work" or "deep embeddings failed." A frozen deep representation was feasible
for current-week visible-corrosion estimation, reaching accuracy broadly
comparable to the handcrafted-feature baseline. Adding metadata to that
representation did not yield a stable improvement in this saved experiment:
it helped during training but not on the final held-out partition, and the
comparison was not confirmed independently of the selection step. The
representation also offered materially less interpretability than the
explicit handcrafted descriptors of the previous phase — a 512-dimensional
embedding does not, by inspection, reveal which visual cues drive a
prediction the way a rust-percentage threshold does. Combined with the
late-stage underestimation problem and the absence of any progress on
structural or remaining-service-life questions, these results motivated a
return to explicit, interpretable image descriptors, alongside a shift
toward posing more demanding questions about hidden structural condition
rather than continuing to refine current-corrosion estimation alone.

\section{Interpretable Corrosion Quantification and Structural-Condition Feasibility}
\label{sec:phase3}

The fourth phase is the project's largest single body of work in this
period and marks a scientific turning point. It replaced the deep embedding
with a broad, explicitly engineered feature representation — ninety-one
descriptors spanning corrosion amount, spatial distribution along the
specimen, surface morphology, colour in three colour spaces (RGB, HSV, and
Lab), and texture — computed by a Python image-processing pipeline. This
pipeline was inspired by the predecessor thesis's manual image-analysis
workflow but is best described as a Python reconstruction and extension of
that workflow's ideas rather than an exact reproduction of it: the original
workflow combined manual GIMP/BIMP editing steps with a separate MATLAB
analysis stage, whereas this phase implements comparable descriptors
programmatically, in some cases extending beyond what the manual workflow
computed.

This phase was also the first to pose two further questions
beyond current-corrosion estimation: whether visible corrosion could help
estimate a specimen's hidden structural condition, and whether a
degradation trajectory and threshold-crossing screening quantity could be
derived from the resulting estimates. Both required confronting the
asymmetry between dense visible-corrosion supervision and sparse structural
supervision established in Section~\ref{sec:structural-supervision}, and
both were evaluated under all three of the grouped evaluation regimes
defined in Section~\ref{sec:eval-regimes} — the first phase in the project
to do so systematically.

\subsection{Visible-corrosion modelling}
\label{sec:phase3-corrosion}

For the two continuous visible-corrosion targets, several candidate model
families were compared (linear regression, RandomForest, histogram-based
gradient boosting, and a small multilayer perceptron); the specific
family selected as best differs by target and by evaluation regime, and is
not repeated model-by-model here. Two categorical severity targets were
also modelled, using the historical \emph{five}-category ordinal scheme in
place at this phase's dataset snapshot (Section~\ref{sec:dataset-audit}) —
these results must not be described as four-class classification, which
belongs to the later, separate dataset and classification effort
(Section~\ref{sec:representation}, strategy~E).

Table~\ref{tab:phase3-results} reports the best mean result for each target
under each evaluation regime. Under a grouped specimen holdout, both
continuous targets are estimated with useful accuracy (total rust: MAE 1.38,
$R^2=0.646$; peak rust: MAE 4.21, $R^2=0.784$). Robustness to campaign
holdout is reasonable — accuracy declines but both targets retain positive
$R^2$ under leave-one-campaign-out. Robustness to treatment holdout is
markedly weaker: mean $R^2$ across treatment folds becomes negative for both
targets, reflecting that individual treatment groups are small
(Section~\ref{sec:loto}) and that at least one held-out treatment fold
produces a strongly negative per-fold $R^2$, pulling the average down. The
five-category classification targets follow a similar pattern: grouped
holdout accuracy is moderate to good (0.83 for total-rust category, though
macro-F1 of 0.573 reveals the imbalance-sensitive weakness that accuracy
alone would hide), and macro-F1 does not exceed roughly 0.47 under either
stricter regime. Read together, visible-corrosion estimation is the most
robustly image-driven layer of this phase: it degrades under distribution
shift, as expected of any model facing genuinely new treatment or campaign
conditions, but it does not collapse.

\subsection{Structural-condition feasibility}
\label{sec:phase3-structural}

The structural-condition models require more careful description, because
several aspects of their design materially affect what the results can be
taken to show. As established in Section~\ref{sec:structural-supervision},
only 48 rows in the dataset carry a structural label, one per specimen.
The feature set used to predict these targets is not image-only: it
combines specimen/exposure metadata, the \emph{manually assigned} visible-corrosion
labels and their categories (the ground-truth annotations, not the previous
subsection's model predictions of them), and a subset of the engineered
image descriptors. The pipeline's own predicted corrosion values are not
included in this feature list, so the damage stage does not chain a learned
corrosion prediction into a learned damage prediction; it instead poses a
single, direct estimation problem — from metadata, manual labels, and image
descriptors together to structural condition — rather than a
strictly image-based inference.

Under a grouped specimen holdout, both structural targets show a modest but
non-trivial relationship to the available features: wire-area loss reaches
MAE 3.90 ($R^2=0.310$), and ultimate load reaches MAE 0.121~kN
($R^2=0.385$) on this phase's own 38/10 grouped-holdout split. These are
this phase's own historical results, computed on its own 792-row dataset
basis and its own feature engineering, and are not compared numerically
against the later terminal-load-focused results discussed in a subsequent
section, which use a different pipeline, feature set, and dataset
alignment.

Under either stricter evaluation regime, performance deteriorates
substantially. Leave-one-campaign-out mean $R^2$ becomes negative for both
targets ($-0.39$ for wire-area loss, $-8.36$ for ultimate load), and
leave-one-treatment-out is similarly negative and, for ultimate load, more
severely so ($-4.57$). Figure~\ref{fig:phase3-robustness} places these
alongside the visible-corrosion results from the previous subsection: the
contrast between the two rows of bars is the central diagnostic finding of
this phase. The available evidence does not support reliable image-based
recovery of hidden structural condition under either stricter regime tested
here. The visible-corrosion layer is markedly more robust overall — both of
its targets retain positive $R^2$ under leave-one-campaign-out — but it is
not uniformly robust: leave-one-treatment-out is less stable and drives
both visible-corrosion targets to negative mean $R^2$ as well
(Table~\ref{tab:phase3-results}), so the two stricter regimes should not be
treated as interchangeable when comparing the two target families.

\begin{table}[htbp]
\centering
\caption{Representative results by target and evaluation regime, interpretable
multi-family feature phase. Model shown is the best-mean-MAE candidate for
that (target, regime) cell. Rust and wire-area-loss targets are in percent;
ultimate load is in kN. All results computed on this phase's own
792-row dataset basis; see Section~\ref{sec:dataset-audit} for how this
differs from the later canonical basis.}
\label{tab:phase3-results}
\begin{tabular}{p{3.6cm}p{2.5cm}p{1.5cm}p{1.2cm}p{1.3cm}}
\toprule
Target & Regime & Model & MAE & $R^2$ \\
\midrule
Total rust (visible) & Grouped holdout & RF & 1.38 & 0.646 \\
Total rust (visible) & LOTO & MLP & 1.53 & $-0.677$ \\
Total rust (visible) & LOCO & RF & 2.10 & 0.582 \\
\midrule
Peak rust (visible) & Grouped holdout & MLP & 4.21 & 0.784 \\
Peak rust (visible) & LOTO & RF & 4.74 & $-0.077$ \\
Peak rust (visible) & LOCO & HGB & 6.21 & 0.595 \\
\midrule
Wire-area loss (structural) & Grouped holdout & RF & 3.90 & 0.310 \\
Wire-area loss (structural) & LOTO & HGB & 10.76 & $-5.481$ \\
Wire-area loss (structural) & LOCO & HGB & 12.28 & $-0.387$ \\
\midrule
Ultimate load (structural) & Grouped holdout & RF & 0.121 & 0.385 \\
Ultimate load (structural) & LOTO & RF & 0.158 & $-4.569$ \\
Ultimate load (structural) & LOCO & ExtraTrees & 0.624 & $-8.356$ \\
\bottomrule
\end{tabular}
\end{table}

\begin{figure}[htbp]
\centering
\includegraphics[width=0.85\textwidth]{fig_4_3_main3_robustness_by_regime.png}
\caption{Best-mean-model $R^2$ by evaluation regime for the two
visible-corrosion targets and the two structural-condition targets,
interpretable multi-family feature phase. Visible-corrosion models are
substantially more robust than the structural models overall, and both
visible-corrosion targets retain positive $R^2$ under leave-one-campaign-out;
leave-one-treatment-out is less stable and produces negative mean $R^2$ for
both visible-corrosion targets as well, while the structural targets
deteriorate far more severely under both stricter regimes. Generated
directly from the phase's own saved regression-metrics tables
(Table~\ref{tab:phase3-results}).}
\label{fig:phase3-robustness}
\end{figure}

Three further diagnostics, reported directly in this phase's own technical
report, explain why. First, correlation between the manually labelled
visible-corrosion measures and wire-area loss is close to zero across the
48 structural rows (total rust: $r=0.021$; peak rust: $r=0.011$) — visible
corrosion, as measured here, is only weakly related to internal wire loss
even before any modelling is attempted. Second, feature-importance analysis
of the full-data ultimate-load model shows its largest contributions coming
from steel-mesh count and campaign indicator variables, not from image
descriptors — the same campaign-linked metadata dominance later quantified
more thoroughly in Section~\ref{sec:confounding}'s discussion of design
confounding. Third, the "final" structural estimator for each target is fit
on all 48 structural rows with no further held-out partition at that step,
and is then applied to every observation in the dataset — including
non-terminal weeks — to produce an estimated structural condition at each
week. These estimates are the model-derived quantities introduced in
Table~\ref{tab:quantity-availability}: they are the output of a model fit
in-sample on the full structural subset, not a repeated physical
measurement, and the degradation curves built from them in later work
inherit this property.

\subsection{What this phase established}
\label{sec:phase3-synthesis}

This phase's evidence supports several specific conclusions, stated here
without extrapolation beyond what was tested. Interpretable image
descriptors can estimate visible corrosion with useful accuracy, including
under a meaningful campaign-holdout stress test. Visible-corrosion
estimation is considerably better supported by the available evidence than
hidden structural-condition estimation. Structural inference in this
dataset is constrained by the combination of sparse supervision (48 rows,
one per specimen), small and unevenly sized treatment groups, and
campaign-linked design variables that dominate the structural models'
feature importances. The structural-condition stage is a mixed-input
estimation problem — metadata, manual corrosion labels, and image
descriptors together — not an image-only inference.

It is also this phase's own technical report, not a subsequent phase, that
first names campaign confounding and the risk of a "downstream label
shortcut" (structural models drawing signal from manually assigned
corrosion labels and design metadata rather than from independently
learned image content) as explicit concerns. Work described in later
sections that quantifies these problems more rigorously and responds to
them with dedicated diagnostics and a reframed research question should
therefore be understood as acting on a problem this phase had already
identified, not as discovering it for the first time.

That reframing followed directly from the findings above. The sparse,
campaign-confounded structural signal recovered here motivated stronger
attention to out-of-sample structural estimation, explicit
image-versus-metadata ablations, robustness-oriented model selection,
feature cleanup, and monotonicity-aware degradation modelling, and
ultimately a reconsidered structural research question targeting one of
the two directly measured terminal endpoints available in the dataset.

\section{Degradation Modelling and Proxy-RUL Feasibility}
\label{sec:phase3-degradation}

The same phase also carried out a feasibility experiment in the derived
screening quantity introduced in Section~\ref{sec:proxy-rul}: for each of the
48 specimens, a degradation curve was fitted to each of four per-week
series — the two directly observed visible-corrosion measures and the two
model-estimated structural measures from Section~\ref{sec:phase3-structural}
— and projected forward to identify when it would cross an engineering
threshold. This section reports what that experiment actually produced, and
is deliberately brief: the exercise is a feasibility check on the
derivation chain itself, not a second opportunity to validate structural
prediction.

Curve fitting used piecewise-linear and exponential families, selected per
specimen and target. Fit quality was good for the two visible-corrosion
series (median $R^2$ of 0.90 and 0.92) and moderate for the two
model-estimated structural series (median $R^2$ of 0.64 and 0.60) — every
curve was fit to the same 15–18 observation weeks per specimen established
in Section~\ref{sec:longitudinal}, so the structural curves are trajectories
through model-estimated values, not through repeated physical
measurements. The final proxy-RUL for a specimen was defined as the
earliest of three threshold-crossing events read off these curves and
projected up to 104 weeks beyond the last observation: a 25\% wire-area-loss
crossing, an ultimate-load crossing at 80\% of a campaign reference load, or
a crossing of a composite health index combining all four series.

Of the 48 specimens, 33 received a finite proxy-RUL estimate within the
104-week projection horizon; the remaining 15 never crossed any threshold
within that horizon and are right-censored. Of the 33 finite estimates,
23 — the large majority — are exactly zero weeks, meaning the relevant
threshold was already crossed at the specimen's last available observation
rather than projected into the future. Figure~\ref{fig:phase3-rul-histogram}
shows this concentration directly. Only the remaining ten specimens carry a
proxy-RUL that is both finite and genuinely forward-looking, ranging up to
55.5 weeks.

\begin{figure}[htbp]
\centering
\includegraphics[width=0.72\textwidth]{fig_4_4_main3_proxy_rul_histogram.png}
\caption{Distribution of finite proxy-RUL estimates (33 of 48 specimens);
the remaining 15 specimens are right-censored (no threshold crossing within
the 104-week projection horizon) and are not shown on this axis. Most
finite estimates are at or near zero weeks. Reproduced from the phase's own
saved output.}
\label{fig:phase3-rul-histogram}
\end{figure}

This outcome illustrates, concretely, why the quantity is a feasibility
result rather than a validated remaining-service-life estimate. Every stage
of the derivation compounds the uncertainty of the stage before it: the
structural curves are fitted to estimates that are themselves weak under
distribution shift (Section~\ref{sec:phase3-structural}), the curve family
and threshold are modelling choices rather than physical constants
(Section~\ref{sec:proxy-rul}), and the 104-week extrapolation horizon
extends well beyond any observed data for most specimens. The pipeline is a
working demonstration that the four-step derivation from repeated observed
and model-estimated condition series to a threshold-crossing screening
quantity can be executed end to end; it does not, and was not designed to,
constitute evidence that the resulting numbers are accurate estimates of
true remaining service life.

\section{Robustness and Diagnostic Refinement}
\label{sec:phase4a}

The next phase of work continued directly from the structural-condition
findings above, using the same underlying dataset structure but a more
rigorous implementation. This phase's contribution should be stated
precisely: campaign confounding, sparse structural
supervision, and the risk of models drawing on manually assigned labels
rather than independently learned image content were all already named, in
those terms, by the previous phase's own report. This phase's contribution
was to quantify how severe these problems were under systematic testing,
audit the feature set for redundancy, and revise the modelling and
selection procedure in response — acting on a diagnosis the project had
already made, not making it for the first time.

A feature-redundancy audit identified thirteen near-constant columns, one
exact duplicate pair, and thirty-six feature pairs with rank correlation of
0.95 or higher within the engineered feature set — including a pairing
already anticipated in Section~\ref{sec:phase3-corrosion}, since one image
descriptor and the total-rust label are correlated at $\rho=0.9996$, a
near-tautological relationship. More informative for the structural
question is where genuine, non-redundant signal was found: the strongest
univariate correlate of wire-area loss is a single image morphology
feature, but only at $r=0.38$; for ultimate load, the five strongest
univariate correlates are not image features at all, but five metadata and
temporal variables — steel-mesh count, chloride concentration, ageing days,
observation week, and terminal week — each correlated with ultimate load at
exactly the same magnitude, $|r|=0.827$, a direct numerical signature of
the fact that these variables are perfectly aligned with campaign identity
on the 48-row structural subset (Section~\ref{sec:confounding}). This
answers, for the ultimate-load target specifically, where the apparently
predictive structural signal actually originates: from experimental design
metadata that happens to be campaign-linked, not from the corrosion images.

The phase's own robustness benchmark, summarised in
Table~\ref{tab:phase4a-robustness} and Figure~\ref{fig:phase4a-collapse},
quantifies the consequence for the two target families side by side —
deliberately, because the two families behave very differently and must
not be described with a single verdict. Visible-corrosion targets remain
stable or even improve in relative terms under leave-one-campaign-out
(surface total rust: 0.77$\times$ its own grouped-holdout MAE; peak rust:
0.96$\times$), with Spearman correlation staying above 0.97 in both cases.
Structural targets do not: ultimate load's leave-one-campaign-out MAE is
3.24 times its grouped-holdout value, with Spearman falling to 0.35, and
wire-area loss's leave-one-campaign-out Spearman turns negative
($-0.089$). The robustness collapse documented here belongs to the
structural stage; it is not a property of the pipeline as a whole, and the
visible-corrosion layer's continued robustness is as much a part of this
phase's evidence as the structural collapse is.

\begin{table}[htbp]
\centering
\caption{Grouped-holdout and leave-one-campaign-out MAE and Spearman
correlation for the best-performing model per target under each individual
evaluation strategy considered on its own, robustness-quantification phase.
This is a per-strategy diagnostic comparison, not necessarily the phase's
final robustness-weighted model selection (see following paragraph). Rust
targets in percent, wire-area loss as a fraction, ultimate load in kN.
Source: the phase's own saved robustness-benchmark table.}
\label{tab:phase4a-robustness}
\begin{tabular}{p{4.0cm}p{2.4cm}p{1.2cm}p{1.2cm}p{1.4cm}}
\toprule
Target & Regime & Model & MAE & Spearman \\
\midrule
Surface total rust (visible) & Grouped holdout & GB & 0.177 & 0.999 \\
Surface total rust (visible) & LOCO & GB & 0.136 & 0.997 \\
\midrule
Peak rust (visible) & Grouped holdout & RF & 1.097 & 0.992 \\
Peak rust (visible) & LOCO & RF & 1.058 & 0.977 \\
\midrule
Wire-area-loss fraction (structural) & Grouped holdout & RF & 0.105 & 0.280 \\
Wire-area-loss fraction (structural) & LOCO & RF & 0.124 & $-0.089$ \\
\midrule
Ultimate load, kN (structural) & Grouped holdout & RF & 0.174 & 0.778 \\
Ultimate load, kN (structural) & LOCO & RF & 0.564 & 0.353 \\
\bottomrule
\end{tabular}
\end{table}

\begin{figure}[htbp]
\centering
\includegraphics[width=0.75\textwidth]{fig_4_5_main4_relative_mae_collapse_heatmap.png}
\caption{Best-model MAE under each evaluation regime, relative to that
target's own grouped-holdout MAE (1.00 by definition). Visible-corrosion
targets (bottom two rows) stay at or below 1.2$\times$ under every regime;
structural targets (top two rows) inflate sharply under leave-one-campaign-out,
most severely for ultimate load (3.24$\times$). Reproduced from the phase's
own saved diagnostic output.}
\label{fig:phase4a-collapse}
\end{figure}

Table~\ref{tab:phase4a-robustness} should be read as a diagnostic
comparison, not as a statement of which model the phase ultimately
selected: it reports, independently for each target, whichever model
performed best under a given evaluation strategy considered on its own.
The phase's final structural-model selection instead uses the
robustness-weighted procedure described below and is recorded separately
in its own saved best-model table. For ultimate load the two coincide:
RandomForest on the metadata-only feature set is both the per-strategy
diagnostic winner shown above and the phase's final robustness-weighted
selection, with identical metrics. For wire-area-loss fraction they differ:
the per-strategy diagnostic winner in Table~\ref{tab:phase4a-robustness} is
RandomForest, but the phase's final robustness-weighted selection is
CatBoost on the metadata-only feature set (grouped-holdout MAE 0.124,
leave-one-treatment-out MAE 0.116, leave-one-campaign-out MAE 0.121,
grouped-holdout Spearman 0.289) — numerically close to, but not identical
with, the diagnostic table's figures for that target. Both comparisons
agree on the point that matters for this section: regardless of which
specific model is chosen, the structural targets' feature importances and
robustness pattern remain metadata-dominant, in contrast to the
visible-corrosion targets.

Two further changes distinguish this phase's methodology from the previous
one. Structural-target model selection moved from picking whichever model
happened to have the best mean score on a single split, to a weighted
combination emphasising leave-one-campaign-out performance most heavily —
a deliberate response to the previous phase's own observation that
grouped-holdout performance alone was an unreliable guide to structural
robustness. Degradation-curve fitting was also revised: the previous
phase's fits were dominated by a flexible, non-parametric curve family
selected for 44 of 48 specimens, which the project's own review noted can
reach near-zero fitting error by interpolation rather than by capturing a
genuine trend; the revised fitting is more evenly split between that family
and a simpler linear one (21 versus 27 specimens), and degradation status
is now reported as an explicit crossed/not-yet-crossed/censored category
breakdown rather than a single ambiguous "weeks remaining" figure. The
phase's own stated conclusion, after these revisions, is worth reporting
directly: the improved structural models remain metadata-dominant, and the
project does not consider this sufficient support for a claim that surface
image features robustly infer hidden structural damage.

\section{Refocused Terminal Ultimate-Load Study}
\label{sec:phase4b}

The robustness findings above motivated a deliberate change in research
question, not merely a further modelling attempt. Rather than continuing to
target a sparse, model-derived, campaign-confounded hidden-damage estimate,
this phase targets one of the two structural quantities in the dataset
that are directly measured at a single, well-defined terminal time point:
ultimate load, selected here in preference to the other (wire-area loss)
for this refocused analysis. The question posed is correspondingly
specific: after controlling for
specimen design and exposure characteristics, does superficial visible
corrosion provide predictive information about terminal ultimate load
beyond what those design characteristics already carry? Because this
question is about incremental value, not standalone accuracy, every
analysis in this phase compares matched feature-set configurations under
the same model-selection and evaluation procedure, rather than reporting a
single "best" pipeline in isolation.

\subsection{Does corrosion add value over metadata, pooled?}

Table~\ref{tab:phase4b-featuresets} reports the pooled, grouped
cross-validation comparison across seven feature-set configurations, all 48
specimens, ten repeated grouped splits. A metadata-only Ridge regression —
eleven features describing specimen design and exposure, no image content
at all — achieves the lowest mean absolute error (0.173~kN) and the
second-highest Spearman correlation (0.758) of any configuration tested.
Adding HSV colour descriptors to the metadata leaves MAE essentially
unchanged (0.173~kN) while marginally improving Spearman (0.771); adding
RGB descriptors instead makes both metrics worse (MAE 0.184~kN, $R^2=0.584$
versus metadata-only's 0.645). Image descriptors used alone, without
metadata, are clearly the weakest configurations tested (RGB-only MAE
0.253~kN, $R^2=0.269$). Read together, the pooled evidence does not show
image descriptors adding meaningful value over metadata alone, and in one
configuration (RGB) adding them measurably hurts performance.

\begin{table}[htbp]
\centering
\caption{Pooled feature-set comparison, terminal ultimate load, grouped
cross-validation (48 specimens, 10 splits). Model shown is the winning
model for that feature set. Source: the phase's own saved comparison
table.}
\label{tab:phase4b-featuresets}
\begin{tabular}{p{4.1cm}p{1.3cm}p{1.4cm}p{1.1cm}p{1.1cm}p{1.3cm}}
\toprule
Feature set & Features & Model & MAE & $R^2$ & Spearman \\
\midrule
Metadata only & 11 & Ridge & 0.173 & 0.645 & 0.758 \\
Metadata + HSV & 47 & Ridge & 0.173 & 0.635 & 0.771 \\
Metadata + RGB + HSV & 88 & CatBoost & 0.183 & 0.607 & 0.753 \\
Metadata + RGB & 52 & Ridge & 0.184 & 0.584 & 0.748 \\
HSV only & 36 & CatBoost & 0.210 & 0.481 & 0.705 \\
RGB + HSV & 77 & CatBoost & 0.234 & 0.411 & 0.668 \\
RGB only & 41 & CatBoost & 0.253 & 0.269 & 0.501 \\
\bottomrule
\end{tabular}
\end{table}

A subset analysis restricted to the 43 specimens whose terminal corrosion
was above a near-zero onset threshold shows a modest, narrower reversal:
there, the combined metadata-plus-RGB-plus-HSV configuration edges ahead of
metadata-only on MAE (0.163 vs.\ 0.173~kN) and $R^2$ (0.660 vs.\ 0.633),
although metadata-only retains a marginally higher Spearman correlation
(0.803 vs.\ 0.791) even in this subset. This is reported as a narrow,
subset-specific observation consistent with image descriptors being more
competitive once a specimen has developed measurable corrosion, not as a
reversal of the pooled conclusion above.

\subsection{What happens within a single mesh family?}

The specimen design matrix in Section~\ref{sec:experimental-context} means
that campaign identity, steel-mesh count, and chloride level co-vary almost
perfectly. Table~\ref{tab:phase4b-featuresets} is therefore a pooled result
across two structurally different specimen designs, and the phase's own
analysis goes on to ask what the same comparison shows within each
mesh family alone. The answer is an important caution against reading the
pooled MAE values too optimistically: within either the 4-mesh or the
7-mesh group individually (24 specimens each), every feature-set
configuration achieves a \emph{negative} mean $R^2$ (4-mesh: $-0.81$ to
$-3.56$; 7-mesh: $-0.30$ to $-4.65$), even for the configurations whose raw
MAE looks numerically similar to the pooled result (0.14–0.26~kN). A
negative $R^2$ means the model performs worse than simply predicting the
group mean for every specimen. Rank correlation is correspondingly weak
(the best Spearman achieved within 4-mesh is only 0.10; within 7-mesh,
0.46). A low mean-absolute-error number within a small, low-variance
subgroup is not, by itself, evidence of a working predictive relationship,
and this dataset provides a direct illustration of that gap between
absolute error and explained variance.

\subsection{Is the pooled corrosion-load correlation genuine?}

The clearest illustration of the same point uses correlation directly
rather than a fitted model. Pooled across both mesh families, terminal
visible corrosion and terminal ultimate load are correlated at Spearman
$\rho=0.462$ ($p=0.0009$, $n=48$) — a relationship that would ordinarily be
read as a moderate, statistically significant positive association.
Stratified by mesh family, this relationship does not merely weaken: within
the 4-mesh group it is statistically indistinguishable from zero
($\rho=-0.046$, $p=0.83$), and within the 7-mesh group it is negative and
borderline significant in the \emph{opposite} direction ($\rho=-0.401$,
$p=0.052$). Figure~\ref{fig:phase4b-mesh-scatter} shows all three panels
together; the phase's own saved output table labels the pooled figure
explicitly as a "confounded" reference for this reason. This pattern is
evidence of substantial confounding by mesh/campaign group structure — the
pooled trend is strongly influenced by between-group differences (the
lower-mesh, higher-load group also happens to have lower corrosion), rather
than demonstrating a direct effect of corrosion on load-bearing capacity.
It does not, on its own, prove that no corrosion-capacity relationship
could exist under other conditions, nor does it fully decompose the pooled
association into its causal components; it shows only that the pooled
correlation in this dataset cannot be used as evidence for one.

\begin{figure}[htbp]
\centering
\includegraphics[width=0.95\textwidth]{fig_4_6_mesh_stratified_corrosion_vs_load.png}
\caption{Terminal surface total rust percentage versus terminal ultimate
load, stratified by steel-mesh family and pooled. The positive pooled
association is not reproduced within either mesh family: it is
approximately null in the 4-mesh group ($\rho\approx-0.05$) and negative in
the 7-mesh group ($\rho\approx-0.40$). Generated directly from the phase's
own saved specimen-summary table.}
\label{fig:phase4b-mesh-scatter}
\end{figure}

\subsection{Robustness under campaign holdout, and a residual-correction attempt}

Under leave-one-campaign-out — the most demanding regime available, and,
per Section~\ref{sec:confounding}, a test of extrapolation across the mesh
and chloride design space rather than ordinary generalisation — every
feature-set configuration collapses to strongly negative $R^2$ (metadata-only:
MAE 0.465~kN, $R^2=-4.64$). This is consistent with, rather than
contradicting, the pooled-versus-stratified finding above: a model trained
entirely within one mesh/campaign design space has no basis for
extrapolating to a design space it has never seen, regardless of which
features it was given. A further attempt to recover some of this
performance by adding a second-stage, mesh-aware residual correction on top
of the metadata-only baseline was tested and \emph{worsened} MAE, RMSE, and
Spearman relative to the baseline alone — a negative result the phase's own
saved comparison records explicitly rather than omitting.

\subsection{Synthesis}

Taken together, the evidence from this phase supports a specific and
bounded conclusion: for terminal ultimate load, pooled across the observed
specimen designs, specimen and exposure metadata is the strongest available
predictor, and visible corrosion descriptors do not provide a reliable
pooled improvement over it — with a narrow exception in a corrosion-onset
subset that does not overturn the pooled result, and with within-mesh and
cross-campaign evaluation both showing that whatever apparent relationship
exists in the pooled data is confounded by, or fails to extrapolate across,
specimen design. This is a materially more defensible scientific claim than
either "corrosion predicts structural capacity" or "corrosion is
irrelevant to structural capacity" — it is a claim about incremental
predictive value, tested directly, under conditions designed to expose
confounding rather than obscure it.

\section{Four-Class Corrosion Classification and Augmentation}
\label{sec:phase5}

The final activity considered in this chapter is a later and structurally
separate effort, addressing a different task from every phase above: direct
image classification into an ordinal visible-corrosion severity category,
using the raw-image deep-classification strategy introduced in
Section~\ref{sec:representation} (strategy~E). It targets the
\emph{current}, four-level severity schema established in
Section~\ref{sec:dataset-audit} (Table~\ref{tab:category-schema}) — not the
five-level historical schema modelled in
Section~\ref{sec:phase3-corrosion} — and is implemented as a standalone
pipeline with no code dependency on any earlier phase.

The current schema is severely imbalanced: of 791 usable images, 658
(83.2\%) fall in the lowest severity class, against 99, 20, and 14 images
in the three more severe classes respectively (Table~\ref{tab:data-modalities}).
This imbalance directly shaped the data-preparation work carried out so
far. A five-recipe augmentation methodology — spanning brightness/contrast
adjustment, saturation and colour-balance adjustment, blur and noise, a
compound recipe, and a mirror-flip-plus-brightness recipe — was applied to
every one of the 791 original images at a uniform six-fold expansion (one
original plus five augmented copies), producing a \emph{global offline
augmentation inventory} of 4{,}746 images (791 originals plus 3{,}955
augmented copies; Table~\ref{tab:phase5-classdist}). Because the expansion
factor is identical across classes, this inventory preserves the original
class proportions exactly.

Specimens were then assigned to training, validation, and test partitions
at the specimen level, consistent with the grouping principle of
Section~\ref{sec:leakage}, with a leakage-safety rule enforced on top of
that assignment: any augmented copy whose original image belongs to a
validation- or test-assigned specimen is excluded from every split, so
that validation and test contain only original, unaugmented images (75
each) and no augmented copy of a held-out specimen ever reaches the
training set. Because validation and test together account for 150 of the
791 original images, their 750 associated augmented copies ($150\times5$)
are removed from the global inventory before splitting. The resulting
manifests therefore contain 3{,}846 training rows (641 original images plus
3{,}205 of their augmented copies), 75 validation rows, and 75 test
rows — 3{,}996 rows used in total, not the full 4{,}746-image inventory.
The training manifest's own class distribution is 3{,}156/540/78/72
(82.1\%/14.0\%/2.0\%/1.9\%), close to but not identical with the global
inventory's proportions, since which original specimens — and therefore
which of their augmented copies — fall into validation or test is not
itself class-balanced. The majority class remains dominant throughout, and
the project's own documentation identifies class-weighting or an
equivalent technique as a remaining requirement for whichever training
procedure is used next.

\begin{table}[htbp]
\centering
\caption{Global offline augmentation inventory (before leakage-safe
exclusion of held-out specimens' augmented copies), four-class severity
schema. Every class is expanded by the same factor, so relative
proportions within this inventory are unchanged; this table does not
represent the final training-manifest composition (3{,}846 rows; see text).
Source: the effort's own dataset-preparation report.}
\label{tab:phase5-classdist}
\begin{tabular}{ccccc}
\toprule
Class (severity) & Original images & Augmented copies & Total & Factor \\
\midrule
1 (lowest) & 658 & 3{,}290 & 3{,}948 & 6$\times$ \\
2 & 99 & 495 & 594 & 6$\times$ \\
3 & 20 & 100 & 120 & 6$\times$ \\
4 (highest) & 14 & 70 & 84 & 6$\times$ \\
\bottomrule
\end{tabular}
\end{table}

A mirror (horizontal) flip is included among the accepted recipes here,
even though the same transformation is explicitly excluded from the
regression work in earlier phases (Section~\ref{sec:representation}). The
distinction is deliberate and physically grounded rather than
inconsistent: total visible-corrosion severity, the target of this task, is
position-invariant — a mirrored image has the same overall severity — whereas
the strip-position and rust-location descriptors used in the regression
phases encode left-right position on the specimen directly, and a flip
would silently invalidate those features' meaning. Several other
augmentations were considered and rejected on physical-plausibility
grounds: label-blending techniques such as MixUp and CutMix were excluded
because the severity classes are discrete and cannot be meaningfully
interpolated or partially overwritten; synthetic image generation was
excluded as infeasible given fewer than twenty training images in the
smallest classes; and elastic deformation, cutout, large rotation, vertical
flip, and heavy colour shift were all excluded as producing image content
inconsistent with how the specimens were actually photographed. Leakage
safety is enforced by a runtime check in the splitting code, which raises
an error if any augmented image is found in the validation or test
partition, rather than relying on documentation alone.

The status of this effort should be stated plainly. Dataset preparation is
complete: the target schema is defined, the class distribution is
characterised, specimen-level splitting is implemented and verified, and
the augmentation pipeline has been run to produce the expanded training
pool described above. No classifier has been trained, compared, or
evaluated on this data at the time of writing — no ResNet-50 or Vision
Transformer fine-tuning, and no accuracy, macro-F1, or confusion-matrix
result exists for this task in the repository. Data preparation for the
four-class classification task is complete; classifier training and
evaluation remain future work.
